\chapter{Clustering Conclusions}
\label{chap:conclusions}

This research set out to evaluate the ``Privacy vs. Utility'' trade-off in the context of unsupervised learning, specifically questioning whether synthetic data generated via the Classification and Regression Trees (CART) method retains sufficient geometric fidelity for accurate clustering. By analyzing over 1,800 datasets across varying distributions, noise levels, and dimensionalities, we have quantified the structural degradation introduced by synthesis and its impact on K-Means and Hierarchical Clustering.

\section{Key Findings}
\label{sec:key-findings}

The empirical results show that the CART synthesis method introduces specific ``blocky'' orthogonal artifacts due to the nature of decision tree splits. While these artifacts are limited in spherical Normal distributions, they become more visible in skewed Gamma distributions, where the smooth tails of clusters are replaced by more angular synthetic boundaries.

Within this setting, K-Means and Hierarchical Clustering display broadly similar behavior when recovering cluster structure from the synthetic data. Both algorithms perform well when the geometry is simple and the clusters are well separated, and both deteriorate when the synthetic geometry becomes more distorted or the design becomes more complex. Rather than identifying a clear winner, the results suggest that the dominant source of degradation is the synthesis process itself.

The structural fidelity metrics reinforce this interpretation. The Gini analysis indicates that K-Means tends to produce more balanced partitions in skewed settings, whereas Hierarchical Clustering allows more imbalance to remain visible. However, this difference does not translate into a consistent overall advantage in cluster identification. In both methods, centroid drift and variance distortion increase in high-dimensional, overlapping, or otherwise difficult scenarios, confirming that the main limitation lies in preserving the multivariate structure during synthesis rather than in one clustering algorithm clearly outperforming the other.

Across the simulated designs, performance improved as cluster separation increased, and it worsened as the number of clusters, the number of variables, or the correlation between variables increased. In other words, greater separation made the synthetic clustering task easier, whereas more complex cluster structure and feature dependence made it harder for the synthetic data to preserve the geometry needed by either clustering method.

\section{Comparison with previous Work}
\label{sec:comparison}

This study extends the foundational work of Xinnuo Chen \cite{chen2025}, which established the baseline utility of synthetic data. While previous studies largely focused on supervised learning tasks, such as classification accuracy, or global statistical properties, this research specifically addressed the unsupervised domain. Unlike general utility metrics that average performance over a dataset, our introduction of structural metrics---specifically the Mean Centroid Distance and Gini Coefficient---allowed for a granular diagnosis of why algorithms fail. We demonstrated that utility is not a binary property of the data, but an interaction between the generation method and the analysis method. More cautiously, the visual evidence in Figure~\ref{fig:tsne_gamma} suggests that the orthogonal artifacts introduced by CART can interfere with the smooth geometric assumptions often associated with K-Means in non-normal settings, but this effect should be interpreted as a plausible mechanism rather than as a definitive incompatibility unique to K-Means.

\section{Limitations}
\label{sec:limitations-conclusions}

The scope of this research implies certain limitations that must be considered when interpreting the results. First, we exclusively evaluated the CART method implemented in the \texttt{synthpop} package. Other synthesis techniques, such as parametric methods or deep learning approaches, may produce different geometric artifacts that could interact differently with clustering algorithms. Second, the comparison was strictly limited to K-Means and Agglomerative Hierarchical Clustering. Density-based algorithms like DBSCAN or model-based clustering methods were not tested, though they may offer distinct advantages in handling synthetic noise. Finally, the use of pure Normal and Gamma distributions provided a controlled experimental setting but may not fully capture the chaotic noise, outliers, and mixed-type variables often found in real-world administrative data.

\section{Generalization}
\label{sec:generalization-conclusions}

Despite the specific algorithms used, the findings generalize to the broader point that clustering utility depends on how well the synthetic data preserves the geometry required by the downstream method. Algorithms based on different geometric assumptions may react differently to CART-induced distortions, but the present results do not support a strong ranking between centroid-based and connectivity-based approaches. This is crucial for data scientists selecting tools for privacy-preserving analytics, as it suggests that validation against the intended clustering workflow is just as critical as the data quality itself.

\section{Future Research Directions}
\label{sec:future-research}

To further resolve the friction between data privacy and analytical utility, future research should focus on alternative generation methods. Evaluating whether Generative Adversarial Networks (GANs) or Variational Autoencoders (VAEs) can generate smooth, non-orthogonal decision boundaries would be a logical next step to see if they preserve the convexity required for K-Means. Additionally, there is significant potential in developing metric-optimized synthesis loss functions that explicitly penalize the distortion of centroid distances or cluster variance during the training of the generator. Finally, creating a standardized ``Unsupervised Utility Score'' based on the structural metrics defined in this study could help automatically flag synthetic datasets that are unsuitable for clustering tasks before they are released to researchers.

\section{Conclusions}
\label{sec:final-conclusions}

A synthesized comparison of the algorithmic robustness under CART generation is presented in Table~\ref{tab:robustness_summary}.

\begin{table}[H]
  \centering
  \caption[Summary of Algorithmic Robustness to CART Artifacts]{Summary of algorithmic robustness to CART-induced geometric artifacts.}
  \label{tab:robustness_summary}
  \renewcommand{\arraystretch}{1.3}
  \begin{tabular}{@{}lp{4.5cm}p{4.5cm}@{}}
    \toprule
    \textbf{Condition} & \textbf{K-Means} & \textbf{Hierarchical (Ward)} \\ \midrule
    \textbf{Normal Data (Spherical)} & Highly Robust & Highly Robust \\
    \textbf{Gamma Data (Skewed)} & Similar degradation patterns & Similar degradation patterns \\
    \textbf{High Dimensionality} & High Centroid Drift & High Centroid Drift \\
    \textbf{Overlapping Clusters} & Poor Detection & Poor Detection \\ \bottomrule
  \end{tabular}
  \source{Compiled by the author}
\end{table}

The promise of synthetic data is that it offers ``privacy by design'' without compromising analytical truth. This research clarifies the boundaries of that promise. We conclude that synthetic data generated via CART is not utility-neutral; it can alter the geometry of the data in ways that affect cluster recovery under both K-Means and Hierarchical Clustering, with no clear evidence here that one of the two methods is systematically superior.

For organizations leveraging synthetic data, the practical recommendation is therefore not to privilege one of these clustering algorithms in advance, but to validate the chosen method against the geometry preserved by the synthetic release. When the original distribution is simple and well separated, both K-Means and Hierarchical Clustering remain viable. When the data are skewed, overlapping, high-dimensional, or highly correlated, both methods should be interpreted more cautiously because the main risk arises from the synthetic distortion itself. By understanding these geometric interactions, researchers can more effectively unlock the value of sensitive data while maintaining rigorous privacy standards.
